\section{Problem Formulation}\label{sec:problem}
We consider a pretrained VLA model adapted to a target robot with action space $\mathcal A_e$, where $e$ identifies the embodiment. A policy maps camera observations, proprioception, and a language instruction to robot commands. Its action module must match the target controller; transferring a visual-language representation does not by itself provide a suitable action mapping.

Let $o_t=(I_t^{1:V},q_t)$ contain the available camera images and proprioception, $\ell$ denote the language instruction, and $h_t$ collect the observations and controller history available at decision time. The policy generates an action chunk $A_t=(a_t,\ldots,a_{t+H-1})$ conditioned on $(h_t,\ell)$; the controller executes a prefix before replanning. The scene configuration $c$ specifies the episode initialization, including object and target poses. It indexes resets and evaluation groups, rather than an additional textual input to the policy. With binary full-task outcome $Y(\tau)$, the evaluation objective is
\begin{equation}
J(\theta)=\mathbb{E}_{(c,\ell)\sim\rho_{\rm full}}
\mathbb{E}_{\tau\sim P_\theta(\cdot\mid c,\ell)}[Y(\tau)].
\end{equation}
We decompose the task into $K$ ordered skills with binary success predicates $Y_k$ and training reset distributions $\rho_k$. A skill such as approaching an object has its own instruction and success criterion; it is not equivalent to completing pick-and-place. In particular, placement training from a held-object start does not measure complete manipulation from an empty gripper. The composed task uses a single final instruction and its own reset distribution $\rho_{\rm full}$.

Dense feedback is used to bootstrap individual skills. Binary predicates then define skill completion, successful experience for imitation, and milestones during full-task RL. Evaluation always measures the appropriate declared success event, independently of the training reward. We also distinguish the exploration policy from the deployment policy.

For the CDPR evaluation, the action $a_t\in[-1,1]^5$ controls Cartesian translation, yaw, and the gripper. Strict full-task success requires grasp, held lift, intentional release, and native placement, without latched carry-slip or wrong-placement events. LIBERO uses its environment goal verifier~\cite{libero}. These task definitions and scene distributions are different, so their success rates are reported separately.
